\section{Introduction}
\label{sec:intro}

The proliferation of computational finance methodologies has fundamentally transformed
investment decision-making processes, with machine learning and artificial intelligence
emerging as dominant paradigms for financial market analysis. Contemporary research has
demonstrated considerable sophistication in predictive modelling, particularly through
the integration of diverse data modalities including technical indicators, market
microstructure data, and alternative data sources such as social media sentiment and
news analytics \citep{li20,zhang22}. Despite these advancements, a critical disconnect
persists between the generation of predictive signals and their systematic translation
into executable portfolio allocation decisions---a gap that remains largely unaddressed
in the current literature.

\subsection{Current State of Predictive Modeling in Finance}

The evolution of financial forecasting methodology has been a continuing movement
towards greater sophistication and the systematic combination of heterogeneous data
modalities. Initial models were largely based on technical analysis, seeking to predict
future price movement using historically derived price indicators and geometric pattern
recognition \citep{lee07}. The subsequent arrival of machine learning pushed the field
through a complete change. The introduction of algorithms such as support vector
machines, ensemble techniques and multilayer perceptrons, in order to increase the
accuracy of market predictions, marked that change \citep{chou18,liu19}. Deep-learning
systems now constitute the modern frontier and have shown a strong ability to capture
the complex nonlinear relationships that exist within high-frequency financial data
\citep{nayak22}.

Another equally significant trend has been the methodical integration of sentiment
analysis into predictive models. Early attempts in this direction used simple
dictionary-based models \citep{loughran11}, but more recent approaches progressively use
transformer-based large language models (LLMs) to extract complex sentiment and thematic
indications from unstructured financial text \citep{jung24,liu24}. Empirical studies such
as those of \citet{mu23} and \citet{bacco24} show that quantitatively well-founded
measurements of investor sentiment can considerably increase the explanatory value of
traditional technically based models. Going beyond sentiment, \citet{onozo24} point out
the emerging capacity of LLMs to support the nowcasting of macroeconomic features, which
extends the range of predictive characteristics that can be quantitatively analysed.

\subsection{The Signal Fusion Paradigm and Its Limitations}

A current trend in modern quantitative finance research is to combine heterogeneous
types of signals in order to arrive at a stronger forecasting model. An example is the
study of \citet{moodi24}, combining technical indicators and sentiment analysis into a
hybrid deep-learning model. In a similar manner, \citet{farimani24} designed a multimodal
adaptive learning framework that dynamically synthesises heterogeneous sources of data to
predict financial markets. Such fused methodologies are consistently shown to perform
better in prediction than their uni-modal counterparts, and are particularly effective in
describing the multifaceted, multi-dimensional drivers of market dynamics.

Despite these substantive methodological developments, an overwhelming deficiency stands
out in the body of extant research: a near-exclusive focus on predictive accuracy, at the
cost of the decision-making mechanisms that are necessary to put a forecast into
practice. With statistical forecasting precision as the primary measure of success, the
predominant notion of success---as \citet{zhang22} note in their systematic review of
decision fusion---is not actual economic value or achieved portfolio performance. The
critical discontinuity that arises out of this singular interest in prediction is what we
refer to herein as the \emph{signal-allocation gap}: the gap between the production of
informative signals and their encoding into a rational process of portfolio construction
that takes account of real-world investment restrictions and investment goals.

\subsection{The Regime Dependency Challenge}

Financial markets are dynamic in nature, evolving over discrete structural regimes,
whether of high volatility, persistent trend or mean reversion, that inherently
condition the behaviour of trading signals and investment strategies \citep{wang19}. An
important gap in the modern literature is the overwhelming oversight of these regime
requirements. Even though the presence of such market states is frequently accepted, its
systematic integration into portfolio-forming structures remains significantly
underdeveloped \citep{kritzman12}.

A number of studies have attempted to deal with market-state variability. As an example,
\citet{huang20} studied dependencies using bicluster mining and fuzzy inference to arrive
at automated trading predictions. Similarly, \citet{alsheebah24} modelled various market
conditions in an emerging economy by combining endogenous and exogenous variables. More
recent work has taken the state variable directly into the allocation decision
\citep{li25,luo26}. The methodologies of the predictive strand, however, are typically
limited to increasing predictive accuracy and do not advance to regime-aware portfolio
optimisation.

Accordingly, the key issue is twofold: it involves not only the determination of the
prevailing market regime, but also the dynamic re-tuning of both the signal-fusion
parameters and the underlying optimisation dynamics in response to the properties of the
diagnosed regime.

\subsection{Portfolio Construction: The Neglected Frontier}

Modern studies show a significant lack of methodologies focused on the derivation of
portfolio allocations from predictive signals. An intelligent method of constructing
portfolios based on the prediction of financial returns, named IntPort, was proposed by
\citet{sami25}; the approach is, however, based on statistical averaging heuristics,
which is insufficiently analytical to constitute a modern optimisation model. In line
with this, most reinforcement-learning applications in automated trading, such as the one
\citet{kabbani22} describe, focus on the trading of individual assets rather than on the
allocation of a portfolio.

This lack of explicit research on portfolio construction is particularly worrying given
the underlying principles of modern portfolio theory, which emphasise the imperative of
the covariance structure and the diversification of risk \citep{markowitz52}. The
existing failure to convert individual asset predictions systematically into a single
allocation method is a significant obstacle to the improvement of risk-adjusted returns.
Moreover, the traditional methodology largely ignores critical practical restrictions,
including transaction costs, position size limits and liquidity provision, which are
inseparable from the implementation of a strategy in live investment settings
\citep{demiguel09}.

\subsection{Our Contributions}

This study bridges the identified methodological gaps through several pivotal
contributions to financial decision science:

\begin{itemize}

\item An explicit regime detection framework employs Gaussian mixture models to isolate
discrete market states from volatility and momentum data, extending the principles of
\citet{wang19}. The implementation integrates this regime intelligence into the
subsequent signal fusion and optimisation stages, and the fitted model is reported with
order selection, transition matrix and run lengths rather than asserted.

\item A decision-theoretic fusion mechanism recalibrates the influence of technical and
sentiment indicators according to their conditional performance within each identified
regime, in place of a static weighting scheme. The paper specifies this mechanism in full
and then tests it: the state-dependent weighting hypothesis is \emph{not} supported by
the evidence, and we report that result rather than the specification that flatters it.

\item The research formulates a regime-contingent optimisation paradigm. This framework
uses a shrinkage covariance estimate \citep{ledoit04}, an explicit turnover penalty and a
regime-dependent risk budget, progressing beyond the static assumptions of conventional
mean-variance analysis. The fused signal is converted to a return metric before it enters
the objective, so that the risk model and the signal are commensurate.

\item A backtesting protocol with walk-forward validation is deployed, featuring
regime-aware performance attribution to delineate the contribution of each strategic
component under varying market environments. Components are separated by ablation at
matched average exposure, which isolates the effect of timing risk from the effect of
simply holding less.

\item An out-of-sample holdout is reported separately. The final 665 sessions of the
sample postdate every specification decision taken in this study, and are evaluated only
once, which addresses the concern that a walk-forward backtest reported in full can still
be tuned in aggregate \citep{harvey15,prado16}.

\item Ultimately, the proposed methodology constructs a transparent, systematic pipeline
to traverse the critical divide between multi-modal signal prediction and executable
allocation, directly countering the signal-allocation gap highlighted by
\citet{zhang22}.
\end{itemize}

Collectively, these contributions furnish both a methodological architecture and an
empirical evaluation of a regime-conscious strategy. The framework improves risk-adjusted
performance over an equal-weighted benchmark on every measure reported, though not
significantly, and does so by timing risk rather than by selecting assets.
Section~\ref{sec:lit} reviews the relevant literature, Section~\ref{sec:method} specifies
the framework, Section~\ref{sec:results} reports the evidence and
Section~\ref{sec:conclusion} concludes.
